\ticket{4}{Логическое программирование: вычисление как доказательство. Базовые вычислительные механизмы языка Пролог. Декларативная и процедурная семантика логических программ}

\begin{examplan}
    \item Объяснить идею логического программирования: программа как теория, вычисление как доказательство.
    \item Описать факты, правила, запросы и дизъюнкты Хорна.
    \item Раскрыть унификацию, SLD-резолюцию, бэктрекинг и отсечение.
    \item Сравнить декларативную и процедурную семантику логических программ.
\end{examplan}

\subsection*{Короткая версия для письменного ответа}

\textbf{Логическое программирование} --- парадигма, в которой программа задается как набор логических утверждений: фактов и правил. Выполнение программы понимается как попытка доказать запрос на основе этих утверждений.

Центральная идея: \textbf{программа есть теория}, а \textbf{вычисление есть конструктивное доказательство цели}. Если цель содержит переменные, результатом доказательства становятся найденные подстановки для этих переменных.

\subsection*{Программа, факты, правила и запросы}

В Prolog программа состоит из клауз, обычно записанных в виде фактов и правил.

\begin{verbatim}
parent(john, mary).
ancestor(X, Y) :- parent(X, Y).

?- ancestor(john, Who).
\end{verbatim}

\begin{itemize}
    \item \textbf{Факт} утверждает безусловную истинность отношения.
    \item \textbf{Правило} имеет голову и тело: голова истинна, если истинны все подцели тела.
    \item \textbf{Запрос} задает цель, которую система должна доказать.
    \item \textbf{Дизъюнкты Хорна} --- ограниченный вид формул логики первого порядка, удобный для автоматического вывода.
\end{itemize}

Если запрос доказан, Prolog возвращает \texttt{true} и значения переменных. Если доказательство не найдено, возвращается \texttt{false}.

\subsection*{Унификация}

\textbf{Унификация} --- сопоставление двух термов с целью сделать их синтаксически одинаковыми путем подстановки значений вместо переменных.

\begin{itemize}
    \item Две одинаковые константы унифицируются.
    \item Переменная может быть связана с термом.
    \item Структуры унифицируются, если имеют одинаковый функтор, одинаковую арность и попарно унифицируемые аргументы.
\end{itemize}

Например, \texttt{f(A,b)} и \texttt{f(c,B)} унифицируются с подстановкой \texttt{A = c}, \texttt{B = b}. Наиболее общий унификатор называют MGU.

\subsection*{SLD-резолюция}

\textbf{SLD-резолюция} --- стратегия вывода, используемая Prolog для доказательства целей.

\begin{enumerate}
    \item Берется исходный запрос.
    \item Из списка целей выбирается подцель, обычно самая левая.
    \item Prolog просматривает факты и правила сверху вниз и ищет голову, унифицируемую с выбранной подцелью.
    \item Если найден факт, подцель удаляется.
    \item Если найдено правило, подцель заменяется телом правила.
    \item Найденная подстановка применяется к оставшимся целям.
    \item Если список целей пуст, доказательство успешно.
\end{enumerate}

Таким образом, логический вывод одновременно является и поиском доказательства, и вычислением значений переменных.

\subsection*{Бэктрекинг и отсечение}

\textbf{Бэктрекинг} --- поиск с возвратом. Если текущий путь вывода не приводит к успеху или пользователь запрашивает следующее решение, Prolog возвращается к последней точке выбора и пробует альтернативный факт или правило.

Порядок правил, фактов и подцелей влияет на эффективность и даже на завершимость программы.

\textbf{Отсечение} \texttt{!} управляет бэктрекингом:
\begin{itemize}
    \item \textbf{зеленое отсечение} используется как оптимизация и не меняет логический смысл программы;
    \item \textbf{красное отсечение} меняет множество возможных решений и делает программу менее декларативной.
\end{itemize}

\subsection*{Декларативная и процедурная семантика}

\textbf{Декларативная семантика} отвечает на вопрос: \textit{что означает программа}. Программа рассматривается как набор логических утверждений, а запрос --- как вопрос о логическом следовании.

\textbf{Процедурная семантика} отвечает на вопрос: \textit{как Prolog выполняет программу}. Она включает выбор целей слева направо, просмотр правил сверху вниз, унификацию, SLD-резолюцию, бэктрекинг и отсечение.

В чистом логическом программировании желательно, чтобы декларативный смысл и процедурное поведение совпадали. На практике они могут расходиться: например, неудачный порядок рекурсивных правил может привести к бесконечному поиску, хотя декларативно решение существует.

\begin{important}
    \item В логическом программировании вычисление трактуется как доказательство цели.
    \item Базовые элементы Prolog: факты, правила, запросы, термы и переменные.
    \item Унификация подбирает значения переменных, SLD-резолюция строит доказательство, бэктрекинг перебирает альтернативы.
    \item Декларативная семантика описывает смысл программы, процедурная --- порядок ее выполнения.
    \item Отсечение \texttt{!} полезно для управления поиском, но может разрушать декларативность.
\end{important}
